%The technological evolution of recent decades has profoundly transformed the business landscape, making IT solutions a crucial factor for the success of enterprises. In this context, Elmec Informatica S.p.A. stands out as a significant player in the Information Technology sector in Italy. Founded in 1971 and headquartered in Varese, Elmec combines extensive industry experience with a strong focus on innovation, offering services and solutions that cover a wide range of business needs, offering advanced IT solutions for medium and large enterprises. 
%
%The company manages four datacenters (including one Tier IV certified), provides Hybrid IT services, digital workspace and Device-as-a-Service solutions, and is a key partner for cybersecurity and regulatory compliance. With over 450 certified technicians, Elmec integrates cloud technologies (in collaboration with AWS and GAIA-X), promotes environmental sustainability, and stands out for its innovation through its Technological Campus and ongoing investments in professional training.
%
%The primary context of this internship revolves around the redesign and enhancement of the existing \textit{Service Request Portal (SRP)}, which is essential for managing client service requests efficiently.
%
%A Service Requests is a formal request submitted by a customer to obtain specific services or actions within their IT infrastructure. These requests can encompass a wide range of activities, from creating new user accounts in the Active Directory system to managing user permissions, resetting passwords, and deploying software updates.
%
%Various challenges have been identified within the current architecture, particularly concerning the dynamic forms used by customers for submitting the service requests, the management of customizable fields, and the integration with external data sources.
%
%The main objectives of this internship are:
%\begin{itemize}
%    \item \textbf{Analysis of the Existing System}: Conduct a comprehensive analysis of the current state of the SRP system, identifying weaknesses and inefficiencies that hinder its performance and user experience;
%    \item \textbf{Architectural redesign}: Propose and implement a complete architectural redesign of the system;
%    \item \textbf{Modular and Scalable Solution}: Design a modular and scalable solution that can accomodate various client-specific configurations;
%    \item \textbf{Migrations}: Plan and execute a full migration of databases and other company services/platforms that use SRP.
%    \item \textbf{Fine-tuned LLMs integration}: Explore opportunities for integrating specifically fine-tuned LLM systems into the SR workflow.    
%\end{itemize}

\chapter*{Introduction}
\addcontentsline{toc}{chapter}{Introduction}

This thesis describes the redesign and restructuring of a legacy system for managing Service Requests for an internship at Elmec Informatica, evolving it from a legacy distributed monolith with lots structural limitations (SRP) into a modern, scalable microservices architecture (ServiceRequestManager, SRM).

The legacy system is characterized by a fragmented architecture that compromises scalability, maintainability, and reliability. Key challenges identified includes a monolithic frontend component exceeding 11,000 lines of code where business logic is hard-coded into the presentation layer; a heavy dependency on the Agent installed on customer Domain Controllers, leading to security risks and operational bottlenecks; the use of unversioned logic stored within Stored Procedures and SSIS jobs, making collaboration and error tracking difficult; and significant data integrity issues within the legacy MSSQL database.

To overcome these obstacles, the new SRM architecture leverages a modern technological stack. The backend is implemented using Django and Django REST Framework, utilizing Celery and Redis for asynchronous task management. The system is containerized with Docker and orchestrated via Kubernetes, ensuring high availability and scalability. Furthermore, the dependency on legacy agents is eliminated by integrating an Ansible-based provisioning tool that allows secure, agentless remote execution.

A major innovation of this work is the implementation of an intelligent recommendation system. As the service catalog expanded to over 300 distinct actions, traditional hierarchical navigation became a source of cognitive load for users. This thesis introduces a Retrieval Augmented Generation (RAG) solution using an on-premise Large Language Model, allowing users to discover services using natural language with transparent reasoning to ensure trust and explainability.

The thesis is organized as follows: Chapter 1 analyzes the existing system and its flaws; Chapter 2 outlines the architectural redesign and requirements; Chapter 3 focuses on the practical implementation including CI/CD pipelines and the three-level configuration hierarchy; Chapter 4 describes the phased migration strategy maintaining business continuity; Chapter 5 details the RAG-based recommendation system.
